\section{Conservative discovery with a first-closure stop}
\label{sec:op3-conservative-local-producer}

The following refinement retains the weaker cubic accuracy bound but
removes the slow native constant mode. Zero teleportation is an internal
discovery device; every requested output still has positive target
teleportation. No global conservative obstacle is assumed to exist.

\begin{lemma}[Conservative prefix bounds; \statusproved]
\label{lem:op3-conservative-prefix-bounds}
Let $x\geq0$ have conservative residual $r=e_v-Lx\geq0$, where
$L=D-A$ uses original degrees on the connected graph. Every edge obeys
$|x_i-x_j|\leq1$. If $S=\{i:x_i>0\}$ is proper, then
$\max_i x_i\leq|S|$.

Run the gap-threshold updates of \cref{thm:op3-native-gap-push} with
$\gamma=1$, native accuracy $e$ and $\lambda=e/2$, stopping at the
first completed push whose positive support is the whole graph, or at
an empty queue. Every processed prefix has
\[
 \operatorname{vol}(S)<2/e,\qquad
 \max_i x_i<2/e,\qquad
 P<8/e^3,
\]
where $P$ counts original degrees with multiplicity over all pushes.
In particular the stopped process terminates in finite work, even if
$\lambda\operatorname{vol}(G)<1$ and the exact conservative obstacle
has no stabilizing solution.
\end{lemma}
\begin{proof}
For a strict superlevel $H$ lying between two unequal endpoint values,
sum the residual identity to obtain
\[
 \sum_{i\in H,\ j\notin H,\ j\sim i}(x_i-x_j)
 =\one_{\{v\in H\}}-\sum_{i\in H}r_i\leq1.
\]
Every summand is positive. Choosing $H$ to separate an arbitrary
unequal edge proves its unit difference bound. Equal edges need no
argument. From a maximizing positive vertex follow a simple path to
the first zero coordinate. The path has at most $|S|$ edges, proving
the proper-support maximum bound. Empty support is immediate.

The local update and queue invariants of the native theorem do not
require a killing term. Residuals stay nonnegative and have total mass
exactly one; every previously pushed vertex retains at least
$\lambda d_i$. Immediately before a new activation, its legal residual
exceeds $e d_i$, so $\lambda\operatorname{vol}(S\cup\{i\})<1$
before any degree-sized row buffer is allocated. This gives the strict
volume bound on every prefix.

All proper prefixes have maximum below $2/e$. If one push first makes
support full, only its previously zero coordinate changes. Its new
value is at most $r_i/d_i\leq1$, while all old coordinates were at
most $n-1$. Since $n\geq2$, the new maximum is at most $n-1<2/e$.
The algorithm stops before another full-support push. Each push raises
one potential by $h>e/2$, and all raises are monotone. Consequently
\[
 (e/2)P<\sum_i d_i x_i
 \leq\operatorname{vol}(S)\max_i x_i<4/e^2
\]
after a nonzero number of pushes. This gives the displayed bound on
$P$; the zero-push case is direct. An infinite run would violate this
uniform prefix bound because every pushed degree is at least one.
\end{proof}

\begin{lemma}[Transfer upward from a conservative proper output;
\statusproved]
\label{lem:op3-conservative-upward-transfer}
Suppose the stopped conservative process terminates with an empty queue
on proper support. If the target parameter satisfies
$\bar\alpha=2\alpha/(1+\alpha)\leq e^2/4$, then its same potential
$x$ has original target ACL accuracy $e$.
\end{lemma}
\begin{proof}
The conservative residual $r_0=e_v-Lx$ lies in $[0,e d]$ by the
queue invariant. The target residual is
\[
 r_\alpha=e_v-(D-(1-\bar\alpha)A)x=r_0-\bar\alpha Ax\leq r_0.
\]
On a positive coordinate, $r_{0,i}\geq(e/2)d_i$ and
$(Ax)_i\leq d_i\max_jx_j<(2/e)d_i$, hence
$r_{\alpha,i}\geq0$. On a zero coordinate it equals
$(e_v)_i+(1-\bar\alpha)\sum_{j\sim i}x_j\geq0$ directly.
This proves both original residual inequalities. This is an upward
transfer using an active residual margin, distinct from the earlier
downward floor wrapper.
\end{proof}

\begin{theorem}[An early-closure local reference; \statusproved]
\label{thm:op3-conservative-local-reference}
For every original graph, point seed, $0<\alpha\leq1$ and
$0<\epsappr<1$, a deterministic local exact-real word algorithm
returns original ACL output with support volume below $2/\epsappr$
in $O(\epsappr^{-3})$ work and total allocated words. It reads each
positive vertex's original row once, of total volume below
$2/\epsappr$, and queries each discovered degree once. It uses
$O(1+\epsappr^{-1})$ live graph-state words. Only when the entire
graph has already become positive does it additionally allocate a
dense matrix, of $O(\epsappr^{-2})$ live words. No supplied support,
numerical oracle, target-alpha logarithm or ambient preprocessing is used.
\end{theorem}
\begin{proof}
Set $e=\epsappr$. If $\bar\alpha\leq e^2/4$, use conservative
discovery. Otherwise use the original target-parameter gap push. In
both cases stop at an empty queue or immediately after the first push
making the active-record count equal the discovered-record count.
Every neighbor of every positive vertex already has a record, since
its row was cached on first activation. Equality of these two counters
therefore says the nonempty positive support is closed; connectedness
makes it the entire graph. Testing this costs constant work per push
and does not query ambient $n$.

If the queue empties on proper support, the conservative case is
certified by \cref{lem:op3-conservative-upward-transfer}, and the
damped case is the original native theorem. The conservative work
bound is \cref{lem:op3-conservative-prefix-bounds}. In the damped
case the mass calculation gives
$P<2/(\bar\alpha e)<8/e^3$. Both cases have original support volume
below $2/e$ on every processed prefix.

On first full closure, solve the \emph{unregularized target} system
\[
 (D-(1-\bar\alpha)A)y=e_v
\]
on the known whole graph. It is positive definite and an M-matrix,
so the unique exact solution is nonnegative and has target residual
zero. No claim that the target regularized obstacle remains fully
positive is needed. All original degrees and rows are already cached.
Build a dense matrix by paid binary label lookup, use in-place
$LDL^T$ elimination with its positive Schur pivots, and perform the two
triangular solves and diagonal divisions. This uses $O(n^3)$ work,
total allocated words at most that work, and $O(n^2)$ live matrix words.
Here $n\leq\operatorname{vol}(G)<2/e$, so these costs fit the stated
bound. The computation is implemented arithmetic, not a dense reference
callback standing in for a faster numerical source.

The stable-record, sorted-array, cached-incidence and intrusive-queue
accounting is exactly the native theorem's: first discovery costs
$O(e^{-2})$, and all repeated updates cost $O(P)$. Adding counter
checks and the possible closed solve preserves the cubic bound.
Emit $\bar\alpha D x$ or $\bar\alpha D y$ in one paid pass.
The positive-support labels are unchanged by a proper return and can
only lie in the already known whole graph on the exact branch.
All claims are in the exact-real word model; arbitrarily small target
teleportation may still cause large rational numerators and denominators.
\end{proof}

\paragraph{Implemented audit and remaining rate.}
\path{conservative_local_producer.py} implements every production step,
including the dense target solve. Its independent validators check
original residuals, conservative strict-superlevel flux, prefix maximum
and monotone-value work bounds, counter-based closure, cached graph access,
and exact zero residual after a closed solve. The regime boundary
$\bar\alpha=e^2/4$ is tested exactly. Source and backend hashes,
parameters, counts and work reservations are in
\path{CONSERVATIVE_LOCAL_PRODUCER_AUDIT.json}.
The native interface now permits zero teleportation only with the
first-full-support stop; its original positive-parameter audit is rerun.
This refinement improves the explicit reference and support constant.
Its cubic worst-case rate remains weaker than the existing nearly
quadratic source bound, and general OP3 remains \statusopen.

The full run checks 2,065 original target certificates, 3,223 exact
prefixes and 985 conservative strict-superlevel identities in 1.209
seconds. It includes 29 implemented closed-system solves, 1,158 pushes,
2,061 degree-weighted updates and three private huge-hub compositions.
Target parameters extend to $2^{-1024}$. At $\epsappr=1/8$ on the
two-vertex graph, this method uses two pushes and one exact solve,
compared with 1,744 pushes by the earlier floor reference. Both satisfy
the same target ACL accuracy, although their internal tolerances differ.
